\section{Données et construction des variables}\label{sec:donnees}

Les sociétés privées proviennent d'Orbis, base de données du vendeur Bureau van Dijk, consultée via
Wharton Research Data Services, pour les comptes annuels de 2013 à 2022, étendus à 2010 pour
l'analyse des deux chocs. Les sociétés cotées proviennent de Compustat Global, du vendeur
S\&P Global Market Intelligence, même période et même accès. Le cadrage au niveau des ménages
provient des comptes de secteur et des comptes financiers d'Eurostat, vendeur Office statistique de
l'Union européenne, en accès ouvert.

L'échantillon de sociétés privées retient la France et cinq contrôles comparables : l'Allemagne,
l'Italie, l'Espagne, la Belgique et les Pays-Bas. Une société entre dans le panel équilibré si elle
déclare un capital social au moins 6 années sur la fenêtre. Une observation par société et par année
est conservée, le dépôt au plus grand total de bilan en cas de doublon. Le panel principal compte
16\,177\,643 années-sociétés portées par 2\,138\,541 sociétés. Les variables continues sont
winsorisées aux 1\textsuperscript{er} et 99\textsuperscript{e} centiles. L'échantillon coté compte
24\,686 années-sociétés sur les sociétés cotées des six pays.

La variable expliquée de l'hypothèse H1 est la \textit{variation du capital}, notée $g_{it}$ pour la
société $i$ l'année $t$, variation annuelle du logarithme du capital souscrit et libéré $K_{it}$,
\begin{equation}\label{eq:var-capital}
  g_{it} = \ln(1 + K_{it}) - \ln(1 + K_{i,t-1}).
\end{equation}
Nous lui adjoignons la \textit{hausse du capital}, indicatrice $a_{it}$ valant un si le capital
souscrit a augmenté d'une année sur l'autre et zéro sinon. Pour l'hypothèse H2, le \textit{levier}
$\ell_{it}$ est le total des dettes $L_{it}$ rapporté au total de l'actif $A_{it}$,
\begin{equation}\label{eq:levier}
  \ell_{it} = L_{it} / A_{it},
\end{equation}
la \textit{croissance de la dette} est la variation du logarithme de la dette financière, dette à
long terme plus emprunts courants, et l'\textit{investissement} est la variation du logarithme des
immobilisations corporelles. Pour la source des fonds et pour l'hypothèse H3, les \textit{réserves
sur actif} $r_{it}$ sont les capitaux propres moins le capital souscrit, rapportés au total de
l'actif,
\begin{equation}\label{eq:reserves}
  r_{it} = (E_{it} - K_{it}) / A_{it},
\end{equation}
où $E_{it}$ est le total des capitaux propres, et le \textit{ratio de fonds propres} $e_{it}$ est
\begin{equation}\label{eq:fonds-propres}
  e_{it} = E_{it} / A_{it}.
\end{equation}
Au niveau coté, l'\textit{émission de capital} est la trésorerie levée par vente d'actions rapportée
à l'actif. Au niveau des ménages, la \textit{distribution} $Y_{ct}$ est le logarithme des dividendes
distribués reçus par les ménages du pays $c$ l'année $t$.

Le modèle ne retient qu'une variable de contrôle au niveau de la société, la \textit{taille}
$s_{it}$, logarithme du total de l'actif,
\begin{equation}\label{eq:taille}
  s_{it} = \ln(1 + A_{it}).
\end{equation}
Les effets fixes sont de deux ordres : un effet fixe de société $\alpha_i$, qui absorbe toute
différence de niveau stable entre sociétés, et un effet fixe d'année $\gamma_t$, qui absorbe tout
choc annuel commun aux six pays. L'âge de la société n'est pas introduit comme contrôle : étant la
différence exacte entre l'année courante et l'année de constitution, il est déjà absorbé par la
combinaison des effets fixes de société et d'année. Les erreurs sont regroupées au niveau du pays,
l'unité à laquelle le traitement est affecté. Les définitions complètes, les champs source et les
formules figurent dans la table~\ref{table:definitions-variables} en annexe, que la
table~\ref{table:stats-descriptives} des statistiques descriptives accompagne.
